
Interpreting biological mechanisms requires identifying the functional roles of molecular components and explaining how their interactions produce a response. A transcription factor (TF) may regulate its target genes by acting at cis-regulatory elements, relying on chromatin regulators, cooperating with other TFs or regulating RNA intermediates~\cite{dong2025znf143,gallant2012deltanp63alpha,biddie2011transcription,yoshioka2012micromanaging}. Evaluating a model's ability to identify TF-to-gene regulatory mechanisms is therefore essential when assessing its suitability for interpreting perturbation experiments. Different regulatory mechanisms can produce the same direction of expression change. Consequently, correctly predicting whether a target gene is upregulated or downregulated does not establish that a model has identified the mechanism responsible for that response.

Existing works evaluated related aspects of biological prediction. BEELINE evaluates gene regulatory network reconstruction~\cite{pratapa2020benchmarking}, GEARS predicts transcriptional responses to multi-gene perturbations~\cite{roohani2024predicting}, PerturbQA asks questions about perturbation experiments~\cite{wu2025contextualizing}.  These tasks assess regulatory relationships or perturbation outcomes, but their prediction targets do not require identifying the specific molecular mechanism underlying an individual TF–gene relationship. PertReason directly evaluates cell-state-conditioned mechanistic explanations using perturbation data and knowledge-graph paths~\cite{kim2026pertreason}. CellChain tests how predictions depend on controlled changes in cross-scale evidence~\cite{xia2026cellchain}. These approaches provide valuable assessments of pathway reconstruction and evidence use. MeBench complements these evaluations by distinguishing experimentally supported modes of TF--gene regulation and testing the required components and their causal roles in both forward and inverse intervention reasoning. 

%Addressing this question requires reference annotations that explicitly connect each regulatory category to the experimental evidence supporting its mediator and causal dependencies. For each regulatory pair in our benchmark, explicit experimental evidence requirements define the reference category, mediator and causal dependencies. Complementary intervention tasks then assess the models’ understanding of these mechanisms.

We assess models’ mechanistic understanding through the content of their answers. Each model must identify the experimentally supported regulatory mechanism, explain the role of the required mediator and the relevant causal relationships, and use these relationships to explain the response to a perturbation. An answer that names the correct molecules but reverses their roles does not meet these criteria. Neither does a correct prediction of response direction without the required mechanistic explanation. We collected evidence from published studies to support this assessment. TF occupancy suggests a potential regulatory interaction, but binding alone does not demonstrate an effect on gene expression~\cite{cusanovich2014functional}. Targeted enhancer perturbations can establish functional links between regulatory elements and their target genes~\cite{fulco2019activity}, while genome-scale Perturb-seq measures expression responses to genetic perturbations~\cite{replogle2022mapping}. Identifying the mechanism behind a response requires further evidence showing how the relevant molecular components contribute to that response and which factors or interactions are necessary~\cite{gallant2012deltanp63alpha}. These distinctions form the basis of the evidence requirements for each regulatory mechanism in our benchmark.

The MeBench evaluates models’ ability to reason about regulatory mechanisms through three tasks. Task 1 asks models to identify how the TF regulates its target gene. Task 2 requires them to apply that mechanism to predict and explain the target gene’s response to a given perturbation. Task 3 requires inverse reasoning: inferring the perturbation type from the response direction and the perturbed entity. A shared five-criterion rubric evaluates the mechanistic reasoning in Tasks 2 and 3. Joint correctness requires both a correct task-specific answer and an explanation that satisfies all five criteria; a correct prediction alone is therefore insufficient for full credit. We evaluated six models using 171 curated TF–to-gene pairs. We assessed their mechanistic understanding from their responses to the given inputs, focusing on their ability to apply biological knowledge to identify regulatory mechanisms and explain perturbation responses. Evaluation results revealed that current LLMs have limitations in identifying regulatory mechanisms of molecular components. In the forward and inverse reasoning tasks, models sometimes identified the relevant molecules but incorrectly characterized functions. Under the benchmark conditions, the evaluated systems do not consistently recover the evidence-supported dependencies connecting regulatory components to target responses.








